Table~\ref{tab:app_main} restates the main robustness result in compact form.

\begin{table}[H]
\centering
\begin{tabular}{lr}
\toprule
Quantity & Value \\
\midrule
Stable strict valid ASR & 30.0\% \\
95\% CI & [14.7\%, 49.4\%] \\
Threshold plateau & 0.65--0.80 \\
Full-repeatability mixed flips & 0 \\
Full-repeatability stable candidates & 120/120 \\
\bottomrule
\end{tabular}
\caption{Main results and robustness summary.}
\label{tab:app_main}
\end{table}

Table~\ref{tab:app_baseline} gives the baseline comparison in expanded form.

\begin{table}[H]
\centering
\begin{tabular}{lrrrrr}
\toprule
Method & n & Mean z & Escape rate & 95\% CI & AUROC \\
\midrule
Baseline pilot 30x5 & 30 & 3.473 & 13.3\% & [3.8\%, 30.7\%] & 0.998 \\
Baseline formal 250x5 & 250 & 3.439 & 10.8\% & [7.2\%, 15.3\%] & 0.985 \\
Proposed stable strict & 30 & -- & 30.0\% & [14.7\%, 49.4\%] & -- \\
\bottomrule
\end{tabular}
\caption{Baseline comparison with uncertainty reporting.}
\label{tab:app_baseline}
\end{table}

Table~\ref{tab:app_ablation} repeats the prompt-component ablation.

\begin{table}[H]
\centering
\begin{tabular}{lrrrr}
\toprule
Method & Mean attack z & Mean z drop & Escape rate & AUROC \\
\midrule
Weak paraphrase & 2.526 & 0.645 & 33.3\% & 0.934 \\
No boundary change & 2.268 & 0.903 & 40.0\% & 0.906 \\
Boundary only & 0.166 & 3.005 & 93.3\% & 0.828 \\
Full E2 & 1.208 & 1.963 & 73.3\% & 0.903 \\
\bottomrule
\end{tabular}
\caption{Prompt-component ablation.}
\label{tab:app_ablation}
\end{table}

Table~\ref{tab:app_failure} summarizes the qualitative failure/quality analysis.

\begin{table}[H]
\centering
\begin{tabular}{lccc}
\toprule
Method & Mean quality & Quality $\geq 4$ & Major-error rate \\
\midrule
Plain paraphrase & 4.43 & 79.4\% & 0.0\% \\
Sentence boundary & 4.24 & 72.7\% & 3.0\% \\
AMR guided & 3.63 & 57.6\% & 9.1\% \\
\bottomrule
\end{tabular}
\caption{Failure and quality summary across rewriting strategies.}
\label{tab:app_failure}
\end{table}

Taken together, these tables show the central tradeoff repeatedly: the strongest raw detector reduction is not the same as the strongest defensible attack. That distinction is the main reason the paper reports both raw escape and strict-valid success.
